% Generated draft: requires mathematical and source review.
\documentclass[11pt]{article}
\usepackage[margin=25mm]{geometry}
\usepackage{fontspec}
\setmainfont{Latin Modern Roman}
\usepackage{amsmath,amssymb,mathtools,unicode-math}
\setmathfont{Latin Modern Math}
\usepackage{enumitem,needspace,xurl,hyperref}
\setlength{\parindent}{0pt}
\setlength{\parskip}{5pt}
\setlength{\emergencystretch}{3em}
\newcommand{\N}{\mathbb{N}}
\newcommand{\Z}{\mathbb{Z}}
\newcommand{\Q}{\mathbb{Q}}
\newcommand{\R}{\mathbb{R}}
\newcommand{\C}{\mathbb{C}}
\newcommand{\F}{\mathbb{F}}
\newcommand{\A}{\mathbb{A}}
\newcommand{\PP}{\mathbb{P}}
\newcommand{\E}{\mathbb{E}}
\newcommand{\eps}{\varepsilon}
\newcommand{\dd}{\,\mathrm{d}}
\begin{document}
\section*{UnsolvedMath ID: 30006471}
\textbf{Source problem code:} catalogue entry
\subsection{Definitions and mathematical statement}
Let \(\mathcal F_0\) denote the polynomial family treated in Theorem 1 of the cited report, and let \(\mathcal E_0\) denote the stratification and multiplicative-character-sum cancellation estimates supplied by that theorem. These symbols are placeholders for material absent from the supplied record; they do not specify a polynomial family or a collection of inequalities.

The source asks for a generalization of Theorem 1 to other polynomial families. The supplied cleaned formulation requests broad, explicitly characterized families of multivariable polynomials for which analogues of the stratification and cancellation estimates in \(\mathcal E_0\) hold. A complete formulation would have to specify the admissible polynomials and characters, the summation domains, the strata and their exceptional cases, the precise estimates, and the parameters in which those estimates are uniform.

The supplied record does not include Theorem 1 or identify \(\mathcal F_0\) and \(\mathcal E_0\). It does not state the coefficient fields, character hypotheses, polynomial nondegeneracy conditions, parameter spaces, stratification bounds, cancellation exponents, or dependence of implied constants. Its original excerpt mentions \(\Sigma(u,v)\) and introduces \(\Lambda\), but the displayed sum is corrupted and its notation is undefined. In particular, the summation ranges, product factors, indices, and exponents cannot be recovered reliably. That expression therefore cannot supply a precise target estimate.

Consequently, the statement above records a research question with unresolved specifications, rather than a fully quantified conjecture. No particular field, polynomial class, or bound is selected here. The supplied record labels the problem open and describes the underlying theorem as proved; those descriptions are supplied source assessments, not independently verified mathematical conclusions.
\subsection{Short English statement}
Find explicitly described families of multivariable polynomials to which the cited Theorem 1's stratification and multiplicative character sum cancellation estimates extend. The theorem, its hypotheses, and its estimates are absent from the supplied record, so a precise quantified conjecture cannot yet be recovered.
\subsection{Sources}
\begin{itemize}
\item Supplied local UnsolvedMath registry record, permanent ID \(30006471\), title \emph{Character Sums for General Polynomial Families}; source problem code \texttt{catalogue entry}. The record belongs to the Oberwolfach Reports Open Problems collection and lists proposed year 2025; no proposer is named.
\item Supplied source citation: \emph{Arithmetic Statistics for Algebraic Objects (2026)}, Oberwolfach Reports, DOI 10.4171/owr/2025/52, pp. 2765--2842. Source URL: \url{https://doi.org/10.4171/owr/2025/52}. The supplied background describes Oberwolfach Reports as records of mathematical workshops and their open-problem discussions. The original excerpt asks for a generalization of Theorem 1 to other polynomial families.
\item Supplied literature reference: \emph{Oberwolfach report 52/2025}, \url{https://doi.org/10.4171/owr/2025/52}. Literature assessment dated exactly \texttt{2026-08-22}: the supplied curation labels the question open and reports that generalization of the proved character-sum theorem to other polynomial families is work in progress. This is supplied dated curation, not an independently verified assessment; its time-sensitive status and cited claims require independent verification.
\item The registry assigns the statement status \texttt{corrected\_verified} and says that the official report was checked, the question isolated from surrounding exposition, and its notation repaired. These are supplied verification claims. The report was not independently consulted for this draft, and the supplied excerpt remains insufficient to reconstruct the theorem or its estimates.
\end{itemize}
\end{document}
